\documentclass{beamer}
% Make sure to run this file with PDFLatex.
%%%%%%%%%%%%%%%%%%%%%%%%%%%%%
% Do not change the next few lines
\usepackage{beamerthemesplit}
\usepackage[english]{babel}

% T1 encoding with Latin Modern provides typewriter braces and a bold
% italic sans font for \emph
\usepackage[T1]{fontenc}
\usepackage{lmodern}

% Stop TOC from appearing at top of every slide
\setbeamertemplate{headline}{}

% Make item symbol square
\setbeamertemplate{itemize item}[square]

% Input my macros
\usepackage{graphicx}

% AMS macros (\text, symbols)
\usepackage{amsmath}
\usepackage{amssymb}

% Package for formatting code segments
\usepackage{verbatim}

% Packages for proof trees
\usepackage{proof}

% context-free grammars
\newcommand{\bnfdef}{\mathrel{\colon\colon\mathord{=}}}
\newcommand{\bnfalt}{\mathrel{\mid}}

% Small typewriter text for inline code; \text makes it work in math
\newcommand{\code}[1]{\text{\small\texttt{#1}}}

% Emphasize code or math
\newcommand{\cemph}[1]{\textcolor{red}{#1}}

\let\OldInfer\infer
\renewcommand{\infer}[3][]
{\OldInfer[\text{\scriptsize\texttt{#1}}]{#2}{#3}}

\newcommand{\cinfer}[3][]{\infer[#1]{\code{#2}}{#3}}
\newcommand{\ttinfer}[3][]{\infer[#1]{\texttt{#2}}{#3}}
\newcommand{\stepto}{\mapsto}


\renewcommand{\emph}[1]{\textbf{\textit{#1}}}

% Now Start entering your information.
%\title{Taming the Wildcards: \\
% Combining Definition- and Use-Site Variance}
% \title{Taming the Wildcards}

\title[Type Theory Tutorial\hspace{2em}\insertframenumber/
\inserttotalframenumber]{Type Theory Tutorial}

\author[John Altidor]{\textbf{John Altidor}}
\date{}
\institute{Logic Seminar Lecture}

% PDF document properties; Beamer sets the title and author itself
\hypersetup{
  pdfsubject={Tutorial on type theory and programming language foundations, using MiniLang},
  pdfkeywords={type theory, type systems, programming languages, static semantics, dynamic semantics, type safety, preservation, progress, substitution lemma, MiniLang}
}

\begin{document}
\frame{\titlepage}

\section{Introduction}

\begin{frame}{Outline}
\begin{itemize}
\item Brief Motivation for Type Systems.

\item Example Type System/Programming Language (PL).
  \begin{itemize}
	\item Presenting $MiniLang$:
	A simple programming language of numbers and strings.
	\item Syntax
	\item Static Semantics (Type Checking)
	\item Dynamic (Operational) Semantics (Evaluation)
	\item Safety Theorems: Preservation + Progress
  \end{itemize}
\item Twelf Tutorial
	\begin{itemize}
	\item Mechanization of $MiniLang$
	\end{itemize}
\end{itemize}
\end{frame}

\section{Motivation}

\begin{frame}{Create Language vs Create Library}
\begin{itemize}
\pause
\item Library Pros:
	\begin{itemize}
	\item Library allows using existing language \emph{infrastructure}

\pause
	\item Smaller learning curve \\
	- Don't need to learn new language constructs.
	\end{itemize}

\pause
\item Library Cons:
  \begin{itemize}
  \item Errors difficult to \emph{detect and debug} w/o a compiler.
		\begin{itemize}
		\item Programs can enter undefined states \\
		(e.g. segmentation fault from reading a non-existing field).
		\end{itemize}

\pause
	\item Requirements \emph{not checked} in the language of the library.
		\begin{itemize}
		\item Leaking confidential information to unauthorized users.
		\end{itemize}
  \end{itemize}
\end{itemize}
\end{frame}

\section{Type Systems}

\begin{frame}{Type Systems}

\begin{itemize}
\item Type System =
\emph{Formally defined} language (calculus) with \emph{types}.

\item
\emph{Types} = Properties/classification over terms (syntax) of a language.

\pause
\item Precisely defining what a language means
  \begin{itemize}
  \item Which programs are allowed in a language?
  \item How does a program execute?
  \item \ldots
  \end{itemize}

\pause  
\item Enables \emph{proving properties} about a language.
  \begin{itemize}
  \item Program is always in a well-defined state throughout
  execution (no segmentation fault).
  \item Can prove properties related to software requirements
  (e.g. information flow).
  \end{itemize}

\pause
\item Compiler \emph{informs} programmers of errors at compile-time.

\pause
\item Best explained with an example: \emph{MiniLang}
\end{itemize}
\end{frame}


\section{Syntax}

\begin{frame}{MiniLang Syntax}
\begin{itemize}
\item \emph{Concrete syntax} is what humans write.
\item \emph{Abstract syntax} is what computers reason over.
\end{itemize}

\[
\begin{array}{llrl@{\hspace{0.5cm}}l}
\mathit{Category}   & \mathit{Item}   & & \mathit{Abstract} & \mathit{Concrete} \\ \\
\mathit{Expression} & e & \bnfdef & x                              & x \\
                    &   & \bnfalt & \code{num[$n$]}                & n \\
                    &   & \bnfalt & \code{str[$s$]}                & \code{'$s$'} \\
                    &   & \bnfalt & \code{+($e_1$; $e_2$)}         & \code{$e_1$ + $e_2$} \\
                    &   & \bnfalt & \code{\^{}($e_1$; $e_2$)}      & \code{$e_1$ \^{} $e_2$} \\
                    &   & \bnfalt & \code{let($x$; $e_1$; $e_2$)}  & \code{let $x$ be $e_1$ in $e_2$}
\end{array}
\]
\end{frame}


\begin{frame}{Example Expressions}
\begin{center}
\resizebox{\linewidth}{!}{%
\begin{tabular}{l|l}
\textit{Abstract} & \textit{Concrete} \\ \hline
\code{+(num[5]; +(num[4]; num[3]))}   & \code{5 + 4 + 3} \\
\\

\code{\^{}(str[john]; \^{}(x; str[doe]))} & \code{'john' \^{} x \^{} 'doe'} \\
\\

\code{let(hours; num[24];}    & \code{let hours be 24 in hours+33} \\
\code{  +(hours; num[33]))}   & \\
\end{tabular}}
\end{center}
\end{frame}

\section{Inference Rules}

\begin{frame}{Inference Rules}
\begin{itemize}
\item Semantics of terms (ASTs) defined w/ inference rules.
\item Rules have the following form.
\end{itemize}

\[
\infer[Rule Label]{\underbrace{J}_{conclusion}}
      {\overbrace{J_1 \quad J_2 \quad \ldots \quad J_n}^{premises}}
\]

\begin{itemize}
\item No premises means axiom.
\end{itemize}
\end{frame}

\section{Type Checking}

\begin{frame}{Static Semantics (Type Checking Rules)}
\[
\cinfer[T.1]{$\Gamma \vdash$ num[$n$]: num}{ }
\qquad
\cinfer[T.2]{$\Gamma \vdash$ str[$s$]: str}{ }
\qquad
\infer[T.3]{\Gamma \vdash x: \tau}{(x, \tau) \in \Gamma}
\]

\[
\cinfer[T.4]{$\Gamma \vdash$ +($e_1$; $e_2$): num}
  {\code{$\Gamma \vdash e_1$: num} & \code{$\Gamma \vdash e_2$: num}}
\quad
\cinfer[T.5]{$\Gamma \vdash$ \^{}($e_1$; $e_2$): str}
  {\code{$\Gamma \vdash e_1$: str} & \code{$\Gamma \vdash e_2$: str}}
\]

\[
\cinfer[T.6]{$\Gamma \vdash$ let($x$; $e_1$; $e_2$): $\tau_2$}
  {\Gamma \vdash e_1: \tau_1
   &
   \Gamma , x : \tau_1 \vdash e_2 : \tau_2
  }
\]
\end{frame}


\begin{frame}{Example Type Derivation}
\begin{center}
\resizebox{\linewidth}{!}{\scriptsize $\displaystyle
\ttinfer[T.6]{$\vdash$ let(hrs; num[24]; +(hrs; num[3])): num}
  {\ttinfer[T.1]{$\vdash$ num[24]: num}{ }
   &
   \ttinfer[T.4]{hrs:num $\vdash$ +(hrs; num[3]): num}
     {\ttinfer[T.3]{hrs:num $\vdash$ hrs: num}{ }
      &
      \ttinfer[T.1]{hrs:num $\vdash$ num[3]: num}{ }
     }
  }
$}
\end{center}
\end{frame}


\begin{frame}{Example Type Check Failure}
\begin{center}
\resizebox{\linewidth}{!}{\scriptsize $\displaystyle
\ttinfer{$\vdash$ let(hrs; num[24]; +(hrs; str[abc]))}
  {\ttinfer[T.1]{$\vdash$ num[24]: num}{ }
   &
   \ttinfer{hrs:num $\vdash$ +(hrs; str[abc]): $Fail$}
     {\ttinfer[T.3]{hrs:num $\vdash$ hrs: num}{ }
      &
      \ttinfer[T.1]{hrs:num $\vdash$ str[abc]: str}{ }
     }
  }
$}
\end{center}
\end{frame}


\section{Evaluation}

\begin{frame}{Dynamic (Operational) Semantics (Evaluation)}
\begin{itemize}
\item Defining how to ``execute'' expressions in $MiniLang$.
\item Specifically, defining a \emph{transition system}/relation
$\stepto$ \\
between expressions to evaluate them to \emph{values}.
\item First, need to define values:
\end{itemize}
\[
\cinfer{num[$n$] value}{ }
\qquad
\cinfer{str[$s$] value}{ }
\]
\end{frame}


\begin{frame}{Dynamic Semantics -- Numerical Addition}
\[
\cinfer[D.1]{+($e_1$; $e_2$) $\stepto$ +($e'_1$; $e_2$)}{e_1 \stepto e'_1}
\]

\[
\cinfer[D.2]{+(num[$n_1$]; $e_2$) $\stepto$ +(num[$n_1$]; $e'_2$)}{e_2 \stepto e'_2}
\]

\[
\cinfer[D.3]{+(num[$n_1$]; num[$n_2$]) $\stepto$ num[$n_3$]}{n_1 + n_2 = n_3}
\]
\end{frame}

\begin{frame}{Dynamic Semantics -- String Concatenation}
\[
\cinfer[D.4]{\^{}($e_1$; $e_2$) $\stepto$ \^{}($e'_1$; $e_2$)}{e_1 \stepto e'_1}
\]

\[
\cinfer[D.5]{\^{}(str[$s_1$]; $e_2$) $\stepto$ \^{}(str[$s_1$]; $e'_2$)}{e_2 \stepto e'_2}
\]

\[
\cinfer[D.6]{\^{}(str[$s_1$]; str[$s_2$]) $\stepto$ str[$s_3$]}{s_1 \code{\^{}} s_2 = s_3}
\]
\end{frame}

\begin{frame}{Dynamic Semantics -- Let Expressions}
\[
\cinfer[D.7]{let($x$; $e_1$; $e_2$) $\stepto$ let($x$; $e'_1$; $e_2$)}{e_1 \stepto e'_1}
\]

\[
\cinfer[D.8]{let($x$; $e_1$; $e_2$) $\stepto$ \cemph{$[e_1 / x]e_2$}}{e_1 \; \code{value}}
\]
\end{frame}

\section{Preservation Theorem}

\begin{frame}{Safety Theorem -- Preservation}
\begin{itemize}
\item If $e$ is a well-typed expression that is not a value,
then performing an evaluation step on $e$ does not change
its type.
\item Formally, if $e : \tau$ and $e \stepto e'$, then $e' : \tau$.

\pause
\item Relates the compile-time analysis (type checking rules)
with the run-time behavior (evaluation rules).
\item Important property for real programming languages.
\end{itemize}
\end{frame}

\section{Preservation Motivation}

\begin{frame}[containsverbatim]{Preservation Motivation}
What if Java did not preserve types during evaluation?
\begin{verbatim}
int x;    // 4 bytes in Java
double y; // 8 bytes in Java
\end{verbatim}
\texttt{x = }$\underbrace{\texttt{x + 8}}_{\text{What if
this evaluated to a double?}}$
\end{frame}

\section{Preservation Proof}

\begin{frame}{Preservation Proof -- Addition Case 1}
Proof by induction on the \emph{possible} typing and
evaluation combinations.

Case: (\texttt{T.4}, \texttt{D.3})
\[
\cinfer[T.4]{+(num[$n_1$]; num[$n_2$]): num}
  {\code{num[$n_1$]: num} & \code{num[$n_2$]: num}}
\]

\[
\cinfer[D.3]{+(num[$n_1$]; num[$n_2$]) $\stepto$ num[$n_3$]}{n_1 + n_2 = n_3}
\]

Using rule \texttt{T.1}:

\[
\cinfer[T.1]{num[$n_3$]: num}{ }
\qed
\]
\end{frame}


\begin{frame}{Preservation Proof -- Addition Case 2}
Case: (\texttt{T.4}, \texttt{D.1})
\[
\cinfer[T.4]{+($e_1$; $e_2$): num}
  {\code{$e_1$: num} & \code{$e_2$: num}}
\qquad
\cinfer[D.1]{+($e_1$; $e_2$) $\stepto$ +($e'_1$; $e_2$)}{e_1 \stepto e'_1}
\]

We assume preservation holds for subexpressions.
Hence, by the \emph{inductive hypothesis},
\texttt{$e_1$: num} and $e_1 \stepto e'_1$ imply
\texttt{$e'_1$: num}.
Rule \texttt{T.4} gives us:

\[
\cinfer[T.4]{+($e'_1$; $e_2$): num}
  {\code{$e'_1$: num} & \code{$e_2$: num}}
\qed
\]
\end{frame}


\begin{frame}{Preservation Proof -- Addition Case 3}

Case: (\texttt{T.4}, \texttt{D.2})
\[
\cinfer[T.4]{+(num[$n_1$]; $e_2$): num}
  {\code{num[$n_1$]: num} & \code{$e_2$: num}}
\]
\[
\cinfer[D.2]{+(num[$n_1$]; $e_2$) $\stepto$ +(num[$n_1$]; $e'_2$)}{e_2 \stepto e'_2}
\]

Since \texttt{$e_2$: num} and $e_2 \stepto e'_2$,
by the inductive hypothesis, \texttt{$e'_2$: num}.

Rule \texttt{T.4} gives us:
\[
\cinfer[T.4]{+(num[$n_1$]; $e'_2$): num}
  {\cinfer[T.1]{num[$n_1$]: num}{ }
   &
   \code{$e'_2$: num}}
\qed
\]
\end{frame}


\begin{frame}{Preservation Proof -- Remaining Cases}
\begin{itemize}
\item Remaining cases in preservation proof apply similar reasoning.
\item We show one more case involving a common lemma.
\end{itemize}
\end{frame}

\section{Substitution}

\begin{frame}{Substitution Lemma}
\begin{itemize}
\item
For a case in the preservation proof, we need the
\emph{Substitution Lemma}:

\item
In words, we can substitute subexpressions that are of the
same type in an expression $e$ without changing the type of $e$.

\pause
\item
Formally: \\
If $\Gamma \vdash e' : \tau'$ and $\Gamma, y : \tau' \vdash e : \tau$,
then $\Gamma \vdash [e' / y]e : \tau$.

\item
Proof of this lemma by induction on the structure of $e$.
\end{itemize}
\end{frame}


\begin{frame}{Preservation Proof -- Let Case}
Case: (\texttt{T.6}, \texttt{D.8})
\[
\cinfer[T.6]{let($x; e_1; e_2$): $\tau_2$}
  {e_1: \tau_1
   &
   x : \tau_1 \vdash e_2: \tau_2
  }
\]
\[
\cinfer[D.8]{let($x$; $e_1$; $e_2$) $\stepto$ $[e_1 / x]e_2$}{e_1 \; \texttt{value}}
\]

Since $e_1: \tau_1$ and $x : \tau_1 \vdash e_2: \tau_2$,
by \emph{substitution lemma}, we have $[e_1 / x]e_2 : \tau_2$.
$\qed$

We have completed the proof of preservation!
\end{frame}


\section{Type Safety}

\begin{frame}{Type Safety}
\begin{itemize}
\item Preservation + Progress = Type Safety
  \begin{itemize}
  \item Progress theorem and proof presented in the paper.
  \end{itemize}

\pause
\item Type safety ensures program behavior is well-defined
throughout its execution.

\pause
\item Proving language properties is important
(e.g. ruling out certain errors, publishing).

\pause
\item But proofs are long, error prone, and difficult
to validate.

\item Automated support for \emph{deriving proofs} and
\emph{checking proofs} of language properties.
	\begin{itemize}
	\item Twelf, Coq, Isabelle, Agda, $\ldots$
	\end{itemize}
\end{itemize}
\end{frame}


\begin{frame}{Summary}
\begin{itemize}
\item Programming languages can be defined using formal
mathematical specification.
	\begin{itemize}
	\item Which programs are allowed.
	\item How a program executes.
	\end{itemize}
\item Formal specification enables proving language properties.
\item Type system = formally defined language with types.
\item Type safety theorems (e.g. preservation) establish
relationship between compile-time analysis and run-time
behavior.
\end{itemize}
\end{frame}
\end{document}
